\chapter{Gaussian 消元与 LU 分解}\label{v0ch:4}
\begin{chaptersummary}
前面已经知道矩阵如何产生以及如何验收近似解。本章从消元建立第一条求解路线：将原系统转为两个三角系统。主元选择处理数值稳定性，分解与回代的分离则解释多右端项的优势。后续稀疏分解和不完全分解都沿用这个机制，只改变存储结构与精确程度。
\end{chaptersummary}
\learninggoals{手算带主元 LU；执行前后代；区分分解、求解和主元增长的成本。}
\paragraph{核心问题}消元为什么能被复用为一个求解算子？

\section{消元、置换与三角系统}
第 $k$ 步以 $a_{kk}$ 为主元，用乘子 $\ell_{ik}=a_{ik}/a_{kk}$ 更新
$a_{ij}\leftarrow a_{ij}-\ell_{ik}a_{kj}$。各步乘子组成单位下三角 $L$，剩余上三角组成 $U$。行置换累计为 $P$，得到 $PA=LU$。
求解依次进行
\[
 y_i=(P\bm b)_i-\sum_{j<i}\ell_{ij}y_j,\qquad
 x_i=\frac{y_i-\sum_{j>i}u_{ij}x_j}{u_{ii}}.
\]
部分主元每列选择剩余行中绝对值最大条目；完全主元同时搜索剩余行列，需额外列置换 $PAQ=LU$，最后恢复未知量顺序。部分主元通常有效，但最坏情形仍可能有较大元素增长。
\section{完整四阶算例}
取
\[
 A=\begin{bmatrix}0&2&1&0\\4&4&0&1\\2&0&3&1\\0&1&1&2\end{bmatrix},
 \quad \bm b=(3,9,6,4)^T.
\]
先交换前两行；第一列的非零乘子为 $\ell_{31}=1/2$。第二列主元为 2，取 $\ell_{32}=-1,\ell_{42}=1/2$；第三列取 $\ell_{43}=1/8$。于是
\[
 L=\begin{bmatrix}1&0&0&0\\0&1&0&0\\1/2&-1&1&0\\0&1/2&1/8&1\end{bmatrix},
 \quad U=\begin{bmatrix}4&4&0&1\\0&2&1&0\\0&0&4&1/2\\0&0&0&31/16\end{bmatrix}.
\]
逐项相乘验证 $LU=PA$。前代得到
$\bm y=(9,3,9/2,31/16)^T$，回代得到 $\bm x=(1,1,1,1)^T$。置换行列式为 $-1$，故 $\det A=-62$。
\conceptbridge{原矩阵可逆，首个主元为什么仍为零？}{可逆性保证存在完整求解，不保证当前行列顺序下每个消元主元非零。行交换重新选择消元坐标；浮点运算中还要避免用很小的主元放大乘子和舍入误差。}{可逆系统 $\longrightarrow$ 合适的消元顺序 $\longrightarrow$ 可稳定执行的三角分解。}
\section{成本、复用与验证}
稠密 LU 分解约 $2N^3/3$ FLOPs，两个三角求解合计约 $2N^2$ FLOPs，原位因子存储 $O(N^2)$。$s$ 个右端的总成本约为分解一次加 $s$ 次回代。块 LU 将大部分更新组织为矩阵乘法；它与逐元素消元数学等价，却具有不同缓存复用。
\costnote{分解因子可用于多右端和迭代改进；矩阵变化通常需要重新分解。报告分解后真实残差，必要时在更高精度计算残差并解修正方程。}
\coderef{sparse_linear_algebra/solvers.py}
\aiexercise{按本章顺序手算四阶矩阵的各步剩余子块、前代和回代，再用库分解交叉核验 $PA=LU$。}{将首个非零主元改为很小的数，比较无主元和部分主元的乘子、元素增长与后向误差。}
\studyroute{略去主元增长因子的最坏界与块算法的浮点误差证明。参考 \cite{golub}。}{partial pivoting; complete pivoting; growth factor; blocked LU; iterative refinement}
